\section{Conclusion}
\label{sec:conclusion}

This paper contributes a frequency-weighted progressive disclosure framework
for instruction documents in multi-agent coding systems and validates it on
\blackrim's 49 KB \claudemd. The conservative trim algorithm (redundancy removal
+ externalisation with pointers) delivers a 21.5\,\% line reduction and 17.1\,\%
cache-amortised token saving without fidelity loss. Beyond the measurement,
we provide an integrated literature review demonstrating why naive learned
compression (trained on retrieval contexts) is mismatched to procedurally-authored
instruction documents, and why knowledge-graph augmentation imposes higher
serialisation cost than benefit.

The broader contribution is methodological. Static-prefix multi-agent coding
systems are an emerging pattern: crews of orchestrated LLM workers, each spawned
with the same large instruction document, operating under prompt-cache amortisation
(80--90\,\% hit rates within a 5-minute window). The cost shape is dominated by
amortisation but the operational burden scales with line count, not amortised tokens.
Maintenance fatigue, new-team onboarding, and fidelity preservation all degrade
with document size. Our work demonstrates that a measurement-first methodology---
frequency analysis, conservative classification, and held-out fidelity gating---
can improve both axes simultaneously, reducing token cost while making the system
simpler to maintain.

Open questions for the community include: (1) building a held-out, adversarial
instruction-following eval suite for static-prefix agents; (2) benchmarking
learned compression specifically on instruction documents to settle whether the
gap from \cref{sec:related} is fundamental or merely unmeasured; and (3) exposing
finer-grained cache-control primitives in LLM client tooling to enable per-section
pinning strategies. This work is the first in a planned two-paper sequence; a
companion submission addresses semantic-similarity routing and cache-sharing
patterns in multi-agent spawning, extending the cost optimisation to dispatch
and scheduling.
